\chapter{Random variables and distributions}
\label{ch:rv}

\begin{question}
How do we turn outcomes into numbers, what information does such a number carry, and how do we describe how likely each value is?
\end{question}

\begin{whatfor}
\begin{itemize}
  \item Defines a random variable as a measurable function, and the $\sigma$-field $\sigma(\xi)$ it generates (its information).
  \item States the Doob--Dynkin lemma: knowing $\xi$ determines $\eta$ exactly when $\eta=f(\xi)$.
  \item Describes a random variable by its distribution, distribution function and density, and lists the distributions used in the book, above all the normal.
\end{itemize}
\skipline{fs:rv}
\end{whatfor}

\begin{prereq}
From Chapter~1: $\sigma$-field, atoms, Borel sets, probability space, continuity of $P$.
\end{prereq}

%=====================================================================
\section{Random variables}

\begin{definition}[Measurable function, random variable]
\label{def:rv}
Let $(\Omega,\F)$ be a sample space with a $\sigma$-field. A function $\xi:\Omega\to\R$ is \dfn{$\F$-measurable} if
\[
  \{\xi\in B\}=\{\omega\in\Omega:\xi(\omega)\in B\}=\xi^{-1}(B)\in\F\qquad\text{for every Borel set }B.
\]
If $(\Omega,\F,P)$ is a probability space, an $\F$-measurable $\xi$ is called a \dfn{random variable} (r.v.).
\end{definition}

Why ask for this? We want $P(\xi\in B)=P(\{\omega:\xi(\omega)\in B\})$ for every Borel set $B$, and $P$ only accepts events in $\F$. Measurability guarantees that every such question is allowed.

\begin{pic}
\begin{tikzpicture}[scale=0.85]
  \draw[thick,defc] (0,0) ellipse (1.6 and 1.2); \node[defc] at (-1.3,1.2) {$\Omega$};
  \draw[thick,thmc,fill=thmc!12] (0.3,-0.1) ellipse (0.8 and 0.55); \node[thmc,font=\scriptsize] at (0.3,-0.1) {$\xi^{-1}(B)$};
  \draw[-{Latex}] (5,-1.5)--(5,1.6) node[above] {$\R$};
  \draw[line width=3pt,thmc] (5,0.2)--(5,0.9); \node[thmc,right] at (5.1,0.55) {$B$};
  \draw[-{Latex},thick] (1.1,0.2) to[bend left=20] node[above] {$\xi$} (4.9,0.55);
\end{tikzpicture}
\picnote{$\xi$ sends each outcome to a number. The event $\{\xi\in B\}$ is the set of outcomes sent into $B$; measurability asks that it be in $\F$.}
\end{pic}

\begin{example}[Number of heads]
\label{ex:heads}
One toss: $\Omega=\{H,T\}$, $\F=\{\varnothing,\{H\},\{T\},\Omega\}$, and $\xi$ = number of heads: $\xi(H)=1$, $\xi(T)=0$. It is enough to check the sets $(-\infty,x]$:
\[
  \xi^{-1}\big((-\infty,x]\big)=\begin{cases}\varnothing,&x<0,\\ \{T\},&0\le x<1,\\ \{H,T\},&x\ge1,\end{cases}
\]
all in $\F$. So $\xi$ is a random variable, and $P(\xi=0)=P(\{T\})$.
\end{example}

Measurability depends on $\F$ as well as on $\xi$. With the smaller $\F=\{\varnothing,\Omega\}$, the same $\xi$ is \emph{not} measurable: $\xi^{-1}((-\infty,\frac12])=\{T\}$ is missing. With that $\F$ you are told nothing, so you cannot answer ``was it tails?''.

%=====================================================================
\section{The information carried by a random variable}

\begin{definition}[$\sigma$-field generated by a random variable]
\label{def:sigma-xi}
The \dfn{$\sigma$-field generated by} $\xi$ is
\[
  \sigma(\xi)=\{\xi^{-1}(B):B\in\B(\R)\}.
\]
It is the smallest $\sigma$-field that makes $\xi$ measurable, and it contains exactly the events that can be decided by looking at the value $\xi(\omega)$. For a family $\xi_i$, $i\in I$, the $\sigma$-field $\sigma(\xi_i,\,i\in I)$ is the smallest one containing every event $\{\xi_i\in B\}$.
\end{definition}

The atoms of $\sigma(\xi)$ are the sets of outcomes that share one value of $\xi$. So if $\xi$ takes $k$ distinct values, $\sigma(\xi)$ has $2^k$ events (\cref{res:count}).

\begin{example}[Three random variables on a die]
\label{ex:sigma-die}
The more $\xi$ reveals, the bigger $\sigma(\xi)$.
\begin{center}
\renewcommand{\arraystretch}{1.25}
\begin{tabular}{@{}lll@{}}
\toprule
$\xi$ & atoms & $\sigma(\xi)$ \\
\midrule
$\xi=7$ (constant) & $\Omega$ & $\{\varnothing,\Omega\}$: $2$ events \\
$\xi=1$ if even, $0$ if odd & $\{2,4,6\},\{1,3,5\}$ & $4$ events \\
$\xi$ = remainder of $\omega$ divided by $3$ & $\{3,6\},\{1,4\},\{2,5\}$ & $8$ events \\
$\xi(\omega)=\omega$ & six singletons & all $64$ subsets \\
\bottomrule
\end{tabular}
\end{center}
\end{example}

\begin{example}[Which outcomes can a family tell apart?]
\label{ex:family}
Two tosses; $\xi_1=1$ if the first toss is $H$, $\xi_2=1$ if the second is $H$, and $S=\xi_1+\xi_2$.
\begin{itemize}
  \item $\sigma(\xi_1)$: atoms $\{HH,HT\}$, $\{TH,TT\}$; $4$ events.
  \item $\sigma(\xi_1,\xi_2)$: every outcome has its own pair $(\xi_1,\xi_2)$, so all four singletons are atoms; $16$ events.
  \item $\sigma(S)$: $S(HT)=S(TH)=1$, so $HT$ and $TH$ cannot be told apart. Atoms $\{HH\},\{HT,TH\},\{TT\}$; $8$ events.
\end{itemize}
Rule: two outcomes are separated by $\sigma(\xi_i,\,i\in I)$ exactly when some $\xi_j$ gives them different values.
\end{example}

%=====================================================================
\section{Functions of a random variable}

\begin{definition}[Borel function]
A function $f:\R\to\R$ is a \dfn{Borel function} if $f^{-1}(B)$ is a Borel set for every Borel set $B$.
\end{definition}

Every continuous function, every piecewise continuous function, and every function you can write with a formula is a Borel function.

\begin{result}[A function of $\xi$ is known once $\xi$ is known]
\label{res:f-xi}
If $f$ is a Borel function and $\xi$ a random variable, then $f(\xi)$ is $\sigma(\xi)$-measurable.
\end{result}
\begin{proof}
Fix a Borel set $B$. Then $\{f(\xi)\in B\}=\{\xi\in f^{-1}(B)\}=\xi^{-1}\big(f^{-1}(B)\big)$. Since $f^{-1}(B)$ is a Borel set, this event is in $\sigma(\xi)$.
\end{proof}

For example, die with $\xi(\omega)=\omega$ and $f(x)=x^2$: $\{f(\xi)\in[10,30]\}=\{\xi\in[\sqrt{10},\sqrt{30}]\cup[-\sqrt{30},-\sqrt{10}]\}=\{4,5\}$.

The converse also holds. It is the precise form of ``$\eta$ is decided by $\xi$''.

\begin{theorem}[Doob--Dynkin lemma]
\label{thm:doob-dynkin}
Let $\xi$ be a random variable. A random variable $\eta$ is $\sigma(\xi)$-measurable if and only if $\eta=f(\xi)$ for some Borel function $f$.
\end{theorem}
\noproof

On a finite $\Omega$ the lemma becomes a test you can run: \emph{$\eta$ is a function of $\xi$ exactly when $\eta$ is constant on every atom of $\xi$.}

\begin{example}[Functions of the number of heads]
Two tosses, $S$ = number of heads, atoms $\{HH\},\{HT,TH\},\{TT\}$.
\begin{itemize}
  \item $\eta=\ind{\text{both }H}$: constant on each atom, and $\eta=f(S)$ with $f(2)=1$, $f(0)=f(1)=0$. Measurable.
  \item $\eta=S^2-1$: $f(x)=x^2-1$. Measurable.
  \item $\eta=\xi_1$ (the first toss): $\xi_1(HT)=1\ne0=\xi_1(TH)$ on the same atom. No $f$ can send the input $1$ to both $1$ and $0$. Not $\sigma(S)$-measurable.
\end{itemize}
Here $\ind{A}$ is the \dfn{indicator} of $A$: it equals $1$ on $A$ and $0$ off $A$.
\end{example}

%=====================================================================
\section{Distribution and distribution function}

\begin{definition}[Distribution; distribution function]
\label{def:df}
The \dfn{distribution} of $\xi$ is the probability measure on Borel sets
\[
  P_\xi(B)=P(\xi\in B),\qquad B\in\B(\R).
\]
Its \dfn{distribution function} (d.f.) is $F_\xi(x)=P(\xi\le x)=P_\xi\big((-\infty,x]\big)$, $x\in\R$.
\end{definition}

The distribution moves the probability from the abstract $\Omega$ onto the real line: each outcome's probability is carried to the point $\xi(\omega)$.

\begin{example}[Heads in two fair tosses]
\label{ex:df-heads}
$\xi$ = number of heads: $P(\xi=0)=\frac14$ ($TT$), $P(\xi=1)=\frac12$ ($HT,TH$), $P(\xi=2)=\frac14$ ($HH$). Adding the mass at values $\le x$,
\[
  F_\xi(x)=\begin{cases}0,&x<0,\\ \frac14,&0\le x<1,\\ \frac34,&1\le x<2,\\ 1,&x\ge2.\end{cases}
\]
\end{example}

\begin{pic}
\begin{tikzpicture}
\begin{axis}[width=8cm,height=4.2cm,xmin=-1,xmax=3,ymin=0,ymax=1.15,xtick={0,1,2},ytick={0.25,0.75,1},yticklabels={$\frac14$,$\frac34$,$1$},xlabel={$x$},ylabel={$F_\xi(x)$}]
  \addplot[defc,very thick,domain=-1:0] {0};
  \addplot[defc,very thick,domain=0:1] {0.25};
  \addplot[defc,very thick,domain=1:2] {0.75};
  \addplot[defc,very thick,domain=2:3] {1};
  \addplot[only marks,mark=*,defc] coordinates {(0,0.25)(1,0.75)(2,1)};
  \addplot[only marks,mark=o,defc,fill=white] coordinates {(0,0)(1,0.25)(2,0.75)};
  \draw[trapc,thick,{Latex}-{Latex}] (axis cs:1.15,0.25)--(axis cs:1.15,0.75) node[midway,right,font=\scriptsize] {jump $=P(\xi=1)=\frac12$};
\end{axis}
\end{tikzpicture}
\picnote{A staircase: flat between values, jumping at each value by that value's probability. The filled dot shows $F$ takes the upper value at a jump (right-continuity).}
\end{pic}

\begin{result}[Properties of a distribution function]
\label{res:df-props}
Every d.f.\ $F=F_\xi$ satisfies:
\begin{enumerate}[label=(\alph*)]
  \item non-decreasing: $x_1<x_2\Rightarrow F(x_1)\le F(x_2)$;
  \item right-continuous: $F(x_n)\to F(x)$ whenever $x_n\downarrow x$;
  \item $\lim_{x\to-\infty}F(x)=0$ and $\lim_{x\to+\infty}F(x)=1$.
\end{enumerate}
Moreover $P(a<\xi\le b)=F(b)-F(a)$, and the jump at $x$ equals the point mass: $F(x)-F(x^-)=P(\xi=x)$, where $F(x^-)$ is the limit from the left.
\end{result}
\begin{proof}
Each property is a statement about nested events.
\begin{steps}
  \item (a) $\{\xi\le x_1\}\subseteq\{\xi\le x_2\}$, so the probabilities are ordered (\cref{res:rules}).
  \item (b) If $x_n\downarrow x$, the events $\{\xi\le x_n\}$ shrink to $\{\xi\le x\}$. Continuity from above (\cref{thm:continuity}) gives $F(x_n)\to F(x)$.
  \item (c) $\{\xi\le -n\}$ shrinks to $\varnothing$ and $\{\xi\le n\}$ grows to $\Omega$; apply continuity again.
  \item $\{\xi\le b\}$ is the disjoint union of $\{\xi\le a\}$ and $\{a<\xi\le b\}$. For the jump, $\{x-\frac1n<\xi\le x\}$ shrinks to $\{\xi=x\}$. \qedhere
\end{steps}
\end{proof}

With a jump, ``$\le$'' and ``$<$'' differ: in \cref{ex:df-heads}, $P(\xi\le1)=\frac34$ but $P(\xi<1)=F(1^-)=\frac14$.

%=====================================================================
\section{Densities and discrete distributions}

\begin{definition}[Density; discrete distribution]
\label{def:density}
A random variable $\xi$ is \dfn{absolutely continuous} with \dfn{density} $f_\xi\ge0$ if
\[
  P(\xi\in B)=\int_B f_\xi(x)\,dx\qquad\text{for every Borel set }B.
\]
It is \dfn{discrete} with values $x_1,x_2,\ldots$ if $P(\xi\in B)=\sum_{x_i\in B}P(\xi=x_i)$ for every Borel set $B$.
\end{definition}

For a density, $F_\xi(x)=\int_{-\infty}^x f_\xi(u)\,du$, so $f_\xi=F_\xi'$ wherever $F_\xi$ is differentiable.

\begin{itemize}
  \item A density is not a probability and can exceed $1$: the uniform density on $[0,\frac12]$ is $2$. Only its integrals are probabilities.
  \item With a density, single points have probability $0$: $P(\xi=x_0)=\int_{x_0}^{x_0}f_\xi=0$. So $P(\xi<b)=P(\xi\le b)$ and endpoints never matter.
\end{itemize}

\begin{definition}[Joint distribution and joint density]
For random variables $\xi_1,\dots,\xi_n$, the \dfn{joint distribution} is $P_{\xi_1,\dots,\xi_n}(B)=P\big((\xi_1,\dots,\xi_n)\in B\big)$ for Borel sets $B\subseteq\R^n$. A \dfn{joint density} $f_{\xi_1,\dots,\xi_n}\ge0$ is a function with
\[
  P\big((\xi_1,\dots,\xi_n)\in B\big)=\int\!\!\cdots\!\!\int_B f_{\xi_1,\dots,\xi_n}(x_1,\dots,x_n)\,dx_1\cdots dx_n .
\]
\end{definition}

\begin{example}[Uniform on the unit square]
\label{ex:square}
$f(x,y)=1$ on $[0,1]^2$: probability is area.
\begin{itemize}
  \item $P(\xi_1\le0.5,\ \xi_2\le0.2)=0.5\times0.2=0.1$.
  \item $P(\xi_1+\xi_2\le1)$ = area of the triangle below the line $x+y=1$ = $\frac12$.
  \item $P(\xi_1=\xi_2)=0$: the diagonal is a line, with zero area.
\end{itemize}
\end{example}

\begin{pic}
\begin{tikzpicture}[scale=3]
  \fill[thmc!20] (0,0)--(1,0)--(0,1)--cycle;
  \fill[defc!35] (0,0) rectangle (0.5,0.2);
  \draw[thick] (0,0) rectangle (1,1);
  \draw[thmc,thick] (1,0)--(0,1) node[above right,font=\scriptsize] {$x+y=1$};
  \draw[trapc,dashed] (0,0)--(1,1) node[right,font=\scriptsize] {$x=y$: area $0$};
  \node[font=\scriptsize] at (0.3,0.5) {area $\frac12$};
  \node[font=\scriptsize,defc] at (0.25,0.1) {$0.1$};
  \node[below,font=\scriptsize] at (1,0) {$1$}; \node[left,font=\scriptsize] at (0,1) {$1$};
\end{tikzpicture}
\picnote{For the uniform density on the square, the probability of a region is its area.}
\end{pic}

%=====================================================================
\section{The distributions used in this book}

\begin{center}
\renewcommand{\arraystretch}{1.35}
\begin{tabular}{@{}lll@{}}
\toprule
Name & Values & $P(\xi=k)$ or density \\
\midrule
Bernoulli$(p)$ & $\{0,1\}$ & $P(\xi=1)=p$, $P(\xi=0)=1-p$ \\
Binomial $B(n,p)$ & $0,1,\dots,n$ & $\binom nk p^k(1-p)^{n-k}$ \\
Uniform $U[a,b]$ & $[a,b]$ & $\frac1{b-a}$ on $[a,b]$ \\
Exponential$(\lambda)$ & $[0,\infty)$ & $\lambda e^{-\lambda x}$ \\
Normal $N(\mu,\sigma^2)$ & $\R$ & $\frac{1}{\sqrt{2\pi\sigma^2}}\,e^{-(x-\mu)^2/(2\sigma^2)}$ \\
\bottomrule
\end{tabular}
\end{center}

The normal distribution is the one we use most.

\begin{definition}[Normal distribution; $\Phi$]
\label{def:normal}
$\xi\sim N(\mu,\sigma^2)$ means $\xi$ has the density $\frac{1}{\sqrt{2\pi\sigma^2}}e^{-(x-\mu)^2/(2\sigma^2)}$. The \dfn{standard normal} is $N(0,1)$; its d.f.\ is written
\[
  \Phi(x)=\int_{-\infty}^x\frac{1}{\sqrt{2\pi}}e^{-u^2/2}\,du .
\]
\end{definition}

\begin{result}[Working with the normal]
\label{res:normal}
Let $Z\sim N(0,1)$.
\begin{enumerate}[label=(\alph*)]
  \item Symmetry: $\Phi(-x)=1-\Phi(x)$; in particular $\Phi(0)=\frac12$.
  \item Scaling: $\mu+\sigma Z\sim N(\mu,\sigma^2)$ for $\sigma>0$. Conversely, if $\xi\sim N(\mu,\sigma^2)$ then $\frac{\xi-\mu}{\sigma}\sim N(0,1)$, so
  \[
    P(\xi\le x)=\Phi\Big(\frac{x-\mu}{\sigma}\Big).
  \]
\end{enumerate}
\end{result}
\begin{proof}
Both are substitutions in the integral.
\begin{steps}
  \item (a) The density $e^{-u^2/2}$ is symmetric about $0$, so the area left of $-x$ equals the area right of $x$, which is $1-\Phi(x)$.
  \item (b) $P(\mu+\sigma Z\le x)=P\big(Z\le\frac{x-\mu}{\sigma}\big)=\Phi\big(\frac{x-\mu}{\sigma}\big)$. Differentiating in $x$ gives $\frac1\sigma\cdot\frac{1}{\sqrt{2\pi}}e^{-(x-\mu)^2/(2\sigma^2)}$, the $N(\mu,\sigma^2)$ density. \qedhere
\end{steps}
\end{proof}

\begin{pic}
\begin{tikzpicture}
\begin{axis}[width=9cm,height=4.2cm,xmin=-3.5,xmax=3.5,ymin=0,ymax=0.45,xtick={-1,0,1},ytick=\empty,xlabel={$x$},samples=100]
  \addplot[draw=none,fill=defc!25,domain=-3.5:-1] {exp(-x^2/2)/sqrt(2*pi)} \closedcycle;
  \addplot[draw=none,fill=thmc!25,domain=1:3.5] {exp(-x^2/2)/sqrt(2*pi)} \closedcycle;
  \addplot[black,thick,domain=-3.5:3.5] {exp(-x^2/2)/sqrt(2*pi)};
  \node[defc,font=\scriptsize] at (axis cs:-2.3,0.12) {$\Phi(-1)$};
  \node[thmc!80!black,font=\scriptsize] at (axis cs:2.4,0.12) {$1-\Phi(1)$};
\end{axis}
\end{tikzpicture}
\picnote{The standard normal density is symmetric, so the two shaded tails are equal: $\Phi(-1)=1-\Phi(1)\approx0.159$.}
\end{pic}

\begin{example}[A normal probability]
$\xi\sim N(2,9)$, so $\mu=2$, $\sigma=3$. Then $P(\xi\le5)=\Phi\big(\frac{5-2}{3}\big)=\Phi(1)\approx0.841$ and $P(\xi>-1)=1-\Phi(-1)=\Phi(1)\approx0.841$.
\end{example}

%=====================================================================
\formulasheet{fs:rv}

\begin{card}{1. Random variable and the information it carries}
\Formula
$\xi$ is a r.v.\ if $\{\xi\in B\}\in\F$ for every Borel $B$ (enough: every $(-\infty,x]$). $\sigma(\xi)=\{\xi^{-1}(B)\}$; atoms = outcomes sharing one value; $k$ values $\Rightarrow 2^k$ events.
\UseWhen
``Is $\xi$ $\F$-measurable?'', ``list $\sigma(\xi)$'', ``how many events in $\sigma(\xi)$?''
\HowTo
\begin{enumerate}
  \item Group the outcomes by the value of $\xi$: these are the atoms.
  \item Measurable w.r.t.\ $\F$ $\iff$ every atom is an event of $\F$.
  \item $\sigma(\xi)$ = all unions of atoms; count $2^k$.
\end{enumerate}
\Caveats
\cav{Measurability depends on $\F$: the same $\xi$ can be measurable for one $\sigma$-field and not for a smaller one.}
\begin{cardex}
Die, $\xi$ = remainder of $\omega$ divided by 3. Atoms $\{3,6\},\{1,4\},\{2,5\}$, so $\sigma(\xi)$ has $2^3=8$ events. Not measurable w.r.t.\ $\{\varnothing,\{1,2,3\},\{4,5,6\},\Omega\}$: the atom $\{1,4\}$ is not in it.
\end{cardex}
\end{card}

\begin{card}{2. Is \texorpdfstring{$\eta$}{eta} a function of \texorpdfstring{$\xi$}{xi}? (Doob--Dynkin)}
\Formula
$\eta$ is $\sigma(\xi)$-measurable $\iff\eta=f(\xi)$ for a Borel $f$. Always: $f(\xi)$ is $\sigma(\xi)$-measurable.
\UseWhen
``Is $\eta$ $\sigma(\xi)$-measurable?'', ``can $\eta$ be written as $f(\xi)$?''
\HowTo
\begin{enumerate}
  \item Find the atoms of $\xi$.
  \item Check $\eta$ is constant on each atom. Yes $\Rightarrow$ $f(\text{value of }\xi)$ = that constant.
\end{enumerate}
\Caveats
\cav{One atom where $\eta$ takes two values kills it.}
\begin{cardex}
Two tosses, $S$ = heads. $\eta=S^2-1$: yes, $f(x)=x^2-1$. $\eta$ = first toss: $HT,TH$ share $S=1$ but differ in $\eta$: \textbf{no}.
\end{cardex}
\end{card}

\begin{card}{3. Distribution function}
\Formula
$F_\xi(x)=P(\xi\le x)$: non-decreasing, right-continuous, $F(-\infty)=0$, $F(\infty)=1$. $P(a<\xi\le b)=F(b)-F(a)$; jump at $x$ = $P(\xi=x)$.
\UseWhen
``Is this a d.f.?'', probabilities of intervals, point masses.
\HowTo
\begin{enumerate}
  \item Check the three properties.
  \item Interval: subtract values of $F$. Point mass: size of the jump.
\end{enumerate}
\Caveats
\cav{At a jump, $P(\xi\le x)\ne P(\xi<x)$. With a density they are equal.}
\begin{cardex}
Heads in two tosses: $F(1)=\frac34$, $F(1^-)=\frac14$, so $P(\xi=1)=\frac12$ and $P(0<\xi\le2)=1-\frac14=\frac34$.
\end{cardex}
\end{card}

\begin{card}{4. Density and joint density}
\Formula
$P(\xi\in B)=\int_Bf_\xi(x)\,dx$, $f_\xi=F_\xi'$. Joint: $P((\xi_1,\xi_2)\in B)=\iint_Bf(x,y)\,dx\,dy$.
\UseWhen
Any probability for a continuous r.v.\ or pair.
\HowTo
\begin{enumerate}
  \item Draw the region $B$.
  \item Integrate the density over it (for a uniform density: area $\times$ height).
\end{enumerate}
\Caveats
\cav{A density can exceed 1; only its integrals are probabilities. Single points and lines have probability 0.}
\begin{cardex}
Uniform on $[0,1]^2$: $P(\xi_1+\xi_2\le1)=\frac12$ (triangle). Exponential(1): $P(\xi>1)=\int_1^\infty e^{-x}dx=e^{-1}\approx0.368$.
\end{cardex}
\end{card}

\begin{card}{5. Normal probabilities}
\Formula
$\xi\sim N(\mu,\sigma^2)$: $P(\xi\le x)=\Phi\big(\frac{x-\mu}{\sigma}\big)$. $\Phi(-x)=1-\Phi(x)$, $\Phi(0)=\frac12$. $\mu+\sigma Z\sim N(\mu,\sigma^2)$.
\UseWhen
Any probability for a normal r.v.\ (later: Brownian motion, Black--Scholes).
\HowTo
\begin{enumerate}
  \item Standardise: subtract $\mu$, divide by $\sigma$.
  \item Read $\Phi$; use symmetry for negative arguments and for ``$>$''.
\end{enumerate}
\Caveats
\cav{Divide by $\sigma$, never by $\sigma^2$.}
\begin{cardex}
$\xi\sim N(2,9)$: $P(\xi\le5)=\Phi(1)\approx0.841$; $P(\xi>-1)=1-\Phi(-1)=\Phi(1)\approx0.841$.
\end{cardex}
\end{card}
